% ============================================================
\chapter{M\'ethodes \`a Noyaux}
\label{ch:noyaux}
\index{m\'ethodes \`a noyaux}
% ============================================================

\noindent\textit{Ce chapitre pr\'esente les m\'ethodes \`a noyaux~: l'astuce
du noyau (\emph{kernel trick}), le th\'eor\`eme de Mercer, les noyaux
classiques, la r\'egression Ridge \`a noyau, l'ACP \`a noyau, le
$k$-means \`a noyau, les espaces de Hilbert \`a noyau reproduisant (RKHS)
et le th\'eor\`eme du repr\'esentant. Nous concluons par les liens avec
les processus gaussiens.}\medskip

% ============================================================
\section{Applications de features et astuce du noyau}
\label{sec:kernel-trick}
\index{kernel trick}\index{astuce du noyau}

\subsection{Motivation~: s\'eparabilit\'e en grande dimension}

\begin{exemple}[XOR non lin\'eairement s\'eparable]
Consid\'erons les points $\{(\pm 1, \pm 1)\}$ dans $\R^2$ avec \'etiquettes
$y = x_1 \cdot x_2$. Aucun hyperplan dans $\R^2$ ne peut les s\'eparer,
mais en ajoutant la feature $\phi_3(\bm{x}) = x_1 x_2$, les donn\'ees
deviennent lin\'eairement s\'eparables dans $\R^3$.
\end{exemple}

\begin{definition}[Application de features]
\label{def:feature-map}
Une \textbf{application de features}\index{application de features}
(ou \emph{feature map}) est une fonction
$\phi : \mathcal{X} \to \mathcal{F}$ qui envoie les donn\'ees d'un
espace d'entr\'ee $\mathcal{X}$ vers un espace de features $\mathcal{F}$
(possiblement de dimension infinie).
\end{definition}

\subsection{D\'erivation formelle de l'astuce du noyau}

De nombreux algorithmes (SVM, Ridge, PCA, $k$-means) ne d\'ependent des
donn\'ees qu'\`a travers des \textbf{produits scalaires}
$\inner{\bm{x}_i}{\bm{x}_j}$. Si l'on travaille dans l'espace des
features, on remplace ces produits par
$\inner{\phi(\bm{x}_i)}{\phi(\bm{x}_j)}_{\mathcal{F}}$, ce qui peut
\^etre co\^uteux si $\dim(\mathcal{F})$ est grand.

\begin{definition}[Fonction noyau]
\label{def:kernel}
Une \textbf{fonction noyau}\index{noyau} est une fonction
$\kappa : \mathcal{X} \times \mathcal{X} \to \R$ telle qu'il existe un
espace de Hilbert $\mathcal{F}$ et une application
$\phi : \mathcal{X} \to \mathcal{F}$ v\'erifiant~:
\begin{equation}
  \kappa(\bm{x}, \bm{x}') = \inner{\phi(\bm{x})}{\phi(\bm{x}')}_{\mathcal{F}}
  \label{eq:kernel-def}
\end{equation}
\end{definition}

\begin{formules}[title=\textbf{L'astuce du noyau}]
Au lieu de~:
\begin{enumerate}
  \item calculer explicitement $\phi(\bm{x}_i)$ pour tout $i$,
  \item former les produits scalaires $\inner{\phi(\bm{x}_i)}{\phi(\bm{x}_j)}$,
\end{enumerate}
on \'evalue directement $\kappa(\bm{x}_i, \bm{x}_j)$, ce qui \'evite de
travailler dans l'espace $\mathcal{F}$ (qui peut \^etre de dimension infinie).

\textbf{Matrice de Gram~:} $\bm{K} \in \R^{n \times n}$ avec
$K_{ij} = \kappa(\bm{x}_i, \bm{x}_j)$.
\end{formules}

\begin{exemple}[Noyau polynomial --- feature map explicite]
Pour $\bm{x}, \bm{x}' \in \R^2$ et $\kappa(\bm{x}, \bm{x}') =
(\bm{x}^\top \bm{x}')^2$~:
\begin{align}
  (\bm{x}^\top \bm{x}')^2 &= (x_1 x_1' + x_2 x_2')^2 \\
  &= x_1^2 {x_1'}^2 + 2 x_1 x_2 x_1' x_2' + x_2^2 {x_2'}^2 \\
  &= \inner{\phi(\bm{x})}{\phi(\bm{x}')}
\end{align}
o\`u $\phi(\bm{x}) = (x_1^2, \sqrt{2}\, x_1 x_2, x_2^2)^\top \in \R^3$.
Le noyau calcule le produit scalaire dans $\R^3$ sans construire $\phi$
explicitement.
\end{exemple}

\subsection{Diagramme~: projection dans l'espace des features}

\begin{center}
\begin{tikzpicture}[
  point/.style={circle, fill, inner sep=1.8pt},
  pospt/.style={point, fill=blue!70},
  negpt/.style={point, fill=red!70}
]
  % Espace d'entree R^2
  \node[font=\bfseries] at (2, 4.2) {Espace d'entr\'ee $\R^2$};
  \draw[->] (0,0) -- (4.2,0) node[right] {\small $x_1$};
  \draw[->] (0,0) -- (0,4) node[above] {\small $x_2$};

  % Points non lineairement separables (cercle interieur vs exterieur)
  \node[pospt] at (1.0, 1.0) {};
  \node[pospt] at (1.2, 1.5) {};
  \node[pospt] at (1.5, 1.2) {};
  \node[pospt] at (1.8, 1.6) {};
  \node[pospt] at (1.4, 1.8) {};

  \node[negpt] at (0.3, 0.5) {};
  \node[negpt] at (3.5, 3.2) {};
  \node[negpt] at (0.5, 3.0) {};
  \node[negpt] at (3.2, 0.4) {};
  \node[negpt] at (3.0, 3.5) {};
  \node[negpt] at (0.4, 3.5) {};

  % Fleche phi
  \draw[-{Stealth}, very thick, purple!70!black] (4.8, 2.0) --
    node[above, font=\small] {$\phi$} (7.2, 2.0);

  % Espace des features F
  \begin{scope}[shift={(8,0)}]
    \node[font=\bfseries] at (2.5, 4.2)
      {Espace des features $\mathcal{F}$};
    \draw[->] (0,0) -- (5,0) node[right] {\small $\phi_1$};
    \draw[->] (0,0) -- (0,4) node[above] {\small $\phi_2$};

    % Points lineairement separables
    \node[pospt] at (1.5, 3.0) {};
    \node[pospt] at (1.8, 2.8) {};
    \node[pospt] at (2.0, 3.2) {};
    \node[pospt] at (2.2, 2.6) {};
    \node[pospt] at (1.7, 3.4) {};

    \node[negpt] at (3.5, 0.8) {};
    \node[negpt] at (4.0, 1.2) {};
    \node[negpt] at (3.8, 0.5) {};
    \node[negpt] at (4.2, 1.5) {};
    \node[negpt] at (3.2, 1.0) {};
    \node[negpt] at (4.5, 0.7) {};

    % Hyperplan separateur
    \draw[thick, green!60!black, dashed] (0.5, 1.0) -- (4.5, 3.8);
    \node[font=\scriptsize, green!50!black] at (4.0, 3.5)
      {S\'eparable~!};
  \end{scope}
\end{tikzpicture}
\end{center}

% ============================================================
\section{Th\'eor\`eme de Mercer et noyaux classiques}
\label{sec:mercer}
\index{th\'eor\`eme de Mercer}

\begin{theoreme}[Mercer]
\label{thm:mercer}
Soit $\kappa : \mathcal{X} \times \mathcal{X} \to \R$ une fonction continue
et sym\'etrique sur un compact $\mathcal{X} \subset \R^d$.
$\kappa$ est un \textbf{noyau d\'efini positif} si et seulement si,
pour tout \'echantillon fini $\{x_1, \ldots, x_n\} \subset \mathcal{X}$,
la matrice de Gram $\bm{K}$ avec $K_{ij} = \kappa(x_i, x_j)$ est
semi-d\'efinie positive~:
\begin{equation}
  \forall \bm{c} \in \R^n, \quad
  \sum_{i=1}^n \sum_{j=1}^n c_i c_j\, \kappa(x_i, x_j) \geq 0
\end{equation}
De plus, $\kappa$ admet une d\'ecomposition~:
\begin{equation}
  \kappa(\bm{x}, \bm{x}') = \sum_{k=1}^{\infty} \lambda_k\,
  \psi_k(\bm{x})\, \psi_k(\bm{x}')
  \label{eq:mercer-decomp}
\end{equation}
o\`u $\lambda_k \geq 0$ sont les valeurs propres et $\{\psi_k\}$ les
fonctions propres de l'op\'erateur int\'egral associ\'e \`a $\kappa$.
\end{theoreme}

\begin{remarque}
La d\'ecomposition de Mercer g\'en\'eralise la d\'ecomposition spectrale
des matrices sym\'etriques d\'efinies positives au cadre fonctionnel.
Elle justifie l'existence de $\phi$ dans~\eqref{eq:kernel-def}~:
on peut poser $\phi(\bm{x}) = (\sqrt{\lambda_1}\,\psi_1(\bm{x}),
\sqrt{\lambda_2}\,\psi_2(\bm{x}), \ldots)$.
\end{remarque}

\subsection{Noyaux classiques}

\begin{formules}[title=\textbf{Noyaux couramment utilis\'es}]
\begin{center}
\renewcommand{\arraystretch}{1.6}
\begin{tabular}{lll}
\toprule
\textbf{Noyau} & \textbf{Expression} & \textbf{Param\`etres} \\
\midrule
Lin\'eaire &
  $\kappa(\bm{x}, \bm{x}') = \bm{x}^\top \bm{x}'$ &
  --- \\
Polynomial &
  $\kappa(\bm{x}, \bm{x}') = (\bm{x}^\top \bm{x}' + c)^p$ &
  $c \geq 0$, $p \in \mathbb{N}^*$ \\
RBF (Gaussien) &
  $\kappa(\bm{x}, \bm{x}') = \exp\!\left(-\dfrac{\norm{\bm{x} - \bm{x}'}^2}{2\sigma^2}\right)$ &
  $\sigma > 0$ \\
Laplacien &
  $\kappa(\bm{x}, \bm{x}') = \exp\!\left(-\dfrac{\norm{\bm{x} - \bm{x}'}_1}{\sigma}\right)$ &
  $\sigma > 0$ \\
\bottomrule
\end{tabular}
\end{center}
\end{formules}

\begin{proposition}[Dimension de l'espace des features du noyau RBF]
\label{prop:rbf-infinite}
Le noyau RBF correspond \`a un espace de features $\mathcal{F}$ de
\textbf{dimension infinie}. En effet, par le d\'eveloppement en s\'erie
de Taylor de l'exponentielle~:
\begin{equation}
  \exp\!\left(-\frac{\norm{\bm{x} - \bm{x}'}^2}{2\sigma^2}\right)
  = e^{-\frac{\norm{\bm{x}}^2}{2\sigma^2}}\,
    e^{-\frac{\norm{\bm{x}'}^2}{2\sigma^2}}\,
    \sum_{k=0}^{\infty} \frac{(\bm{x}^\top \bm{x}')^k}
    {k!\, \sigma^{2k}}
\end{equation}
Chaque terme $(\bm{x}^\top \bm{x}')^k$ engendre un espace de dimension
$\binom{d + k - 1}{k}$, et la somme infinie produit un espace de
dimension infinie.
\end{proposition}

\begin{attention}[title=\textbf{Choix du noyau}]
\begin{itemize}
  \item Le \textbf{noyau lin\'eaire} est adapt\'e lorsque $d \gg n$
    (texte, g\'enomique).
  \item Le \textbf{noyau RBF} est un bon choix par d\'efaut~; le
    param\`etre $\sigma$ (ou $\gamma = 1/(2\sigma^2)$) contr\^ole la
    <<~port\'ee~>> du noyau.
  \item Le \textbf{noyau polynomial} capture les interactions entre
    features jusqu'\`a l'ordre $p$.
\end{itemize}
\end{attention}

% ============================================================
\section{R\'egression Ridge \`a noyau}
\label{sec:kernel-ridge}
\index{r\'egression Ridge!noyau}

\begin{definition}[Kernel Ridge Regression]
\label{def:krr}
La \textbf{r\'egression Ridge \`a noyau} r\'esout~:
\begin{equation}
  \min_{f \in \mathcal{H}_\kappa} \frac{1}{n}
  \sum_{i=1}^n (y_i - f(\bm{x}_i))^2 + \lambda \norm{f}^2_{\mathcal{H}_\kappa}
  \label{eq:krr-functional}
\end{equation}
Par le th\'eor\`eme du repr\'esentant (voir \S\ref{sec:representer}),
la solution s'\'ecrit $f^*(\bm{x}) = \sum_{i=1}^n \alpha_i\,
\kappa(\bm{x}_i, \bm{x})$ avec~:
\begin{equation}
  \bm{\alpha}^* = (\bm{K} + \lambda n\, \bm{I})^{-1} \bm{y}
  \label{eq:krr-alpha}
\end{equation}
o\`u $\bm{K}$ est la matrice de Gram, $K_{ij} = \kappa(\bm{x}_i, \bm{x}_j)$.
\end{definition}

\begin{remarque}
Le co\^ut de calcul est $O(n^3)$ pour l'inversion matricielle et $O(n^2 d)$
pour le calcul de $\bm{K}$. Pour $n$ grand, des approximations comme
Nystr\"om ou les features al\'eatoires de Rahimi-Recht sont utilis\'ees.
\end{remarque}

% ============================================================
\section{ACP \`a noyau (Kernel PCA)}
\label{sec:kernel-pca}
\index{ACP!noyau}\index{kernel PCA}

\begin{definition}[Kernel PCA]
L'\textbf{ACP \`a noyau} effectue une analyse en composantes principales
dans l'espace des features $\mathcal{F}$. Les composantes principales
sont les vecteurs propres de la matrice de covariance dans $\mathcal{F}$~:
\begin{equation}
  \bm{C}_\phi = \frac{1}{n} \sum_{i=1}^n \phi(\bm{x}_i)\,
  \phi(\bm{x}_i)^\top
\end{equation}
\end{definition}

\begin{proposition}[Calcul via la matrice de Gram]
Les projections sur les composantes principales dans $\mathcal{F}$ se
calculent \`a partir de la d\'ecomposition spectrale de la matrice de
Gram centr\'ee $\tilde{\bm{K}}$~:
\begin{equation}
  \tilde{\bm{K}} = \bm{K} - \bm{1}_n \bm{K} - \bm{K} \bm{1}_n
  + \bm{1}_n \bm{K} \bm{1}_n
\end{equation}
o\`u $\bm{1}_n = \frac{1}{n}\bm{1}\bm{1}^\top$. Si $\tilde{\bm{K}}
\bm{v}_k = n \mu_k \bm{v}_k$, la projection du point $\bm{x}$ sur
la $k$-i\`eme composante est~:
\begin{equation}
  z_k(\bm{x}) = \sum_{i=1}^n (v_k)_i\,
  \kappa(\bm{x}_i, \bm{x})
\end{equation}
(apr\`es centrage dans $\mathcal{F}$).
\end{proposition}

% ============================================================
\section{$k$-means \`a noyau}
\label{sec:kernel-kmeans}
\index{k-means!noyau}

\begin{definition}[Kernel $k$-means]
Le \textbf{$k$-means \`a noyau} minimise~:
\begin{equation}
  J = \sum_{c=1}^C \sum_{i \in \mathcal{C}_c}
  \norm{\phi(\bm{x}_i) - \bm{\mu}_c^{\phi}}^2_{\mathcal{F}}
\end{equation}
o\`u $\bm{\mu}_c^{\phi} = \frac{1}{|\mathcal{C}_c|}
\sum_{j \in \mathcal{C}_c} \phi(\bm{x}_j)$ est le centro\"ide du
cluster $c$ dans $\mathcal{F}$.
\end{definition}

\begin{formules}[title=\textbf{Calcul des distances dans $\mathcal{F}$}]
La distance au carr\'e entre $\phi(\bm{x}_i)$ et $\bm{\mu}_c^{\phi}$
s'exprime uniquement en termes de noyaux~:
\begin{equation}
  \norm{\phi(\bm{x}_i) - \bm{\mu}_c^{\phi}}^2 =
  \kappa(\bm{x}_i, \bm{x}_i) -
  \frac{2}{|\mathcal{C}_c|} \sum_{j \in \mathcal{C}_c}
  \kappa(\bm{x}_i, \bm{x}_j) +
  \frac{1}{|\mathcal{C}_c|^2} \sum_{j, l \in \mathcal{C}_c}
  \kappa(\bm{x}_j, \bm{x}_l)
\end{equation}
L'algorithme it\`ere affectation et mise \`a jour comme le $k$-means
classique, mais dans $\mathcal{F}$.
\end{formules}

% ============================================================
\section{Espaces de Hilbert \`a noyau reproduisant (RKHS)}
\label{sec:rkhs}
\index{RKHS}\index{espace de Hilbert \`a noyau reproduisant}

\subsection{D\'efinition et intuition}

\begin{definition}[RKHS]
\label{def:rkhs}
Un \textbf{espace de Hilbert \`a noyau reproduisant} $\mathcal{H}_\kappa$
associ\'e \`a un noyau d\'efini positif $\kappa$ est un espace de Hilbert
de fonctions $f : \mathcal{X} \to \R$ satisfaisant~:
\begin{enumerate}
  \item \textbf{Feature map canonique~:} pour tout $\bm{x} \in \mathcal{X}$,
    la fonction $\kappa(\bm{x}, \cdot) \in \mathcal{H}_\kappa$.
  \item \textbf{Propri\'et\'e de reproduction~:} pour tout
    $f \in \mathcal{H}_\kappa$ et tout $\bm{x} \in \mathcal{X}$~:
    \begin{equation}
      f(\bm{x}) = \inner{f}{\kappa(\bm{x}, \cdot)}_{\mathcal{H}_\kappa}
      \label{eq:reproducing}
    \end{equation}
\end{enumerate}
En particulier,
$\kappa(\bm{x}, \bm{x}') = \inner{\kappa(\bm{x}, \cdot)}{\kappa(\bm{x}', \cdot)}_{\mathcal{H}_\kappa}$.
\end{definition}

\begin{intuition}[title=\textbf{Interpr\'etation du RKHS}]
\begin{itemize}
  \item Le RKHS est l'espace de \emph{toutes les fonctions} que l'on peut
    construire comme combinaisons lin\'eaires (et limites) de
    $\kappa(\bm{x}, \cdot)$, appel\'ees <<~features canoniques~>>.
  \item La norme $\norm{f}_{\mathcal{H}_\kappa}$ mesure la \textbf{r\'egularit\'e}
    de $f$~: les fonctions \`a petite norme sont <<~lisses~>> au sens du
    noyau.
  \item \'Evaluer $f$ en un point $\bm{x}$ revient \`a projeter $f$ sur
    $\kappa(\bm{x}, \cdot)$~: c'est la propri\'et\'e de reproduction.
\end{itemize}
\end{intuition}

\subsection{Construction formelle}

\begin{proposition}[Construction du RKHS]
\label{prop:rkhs-construction}
Le RKHS $\mathcal{H}_\kappa$ est la compl\'etion de l'espace
pr\'e-hilbertien~:
\begin{equation}
  \mathcal{H}_0 = \left\{f = \sum_{i=1}^m \alpha_i\,
  \kappa(\bm{x}_i, \cdot) : m \in \mathbb{N},\;
  \alpha_i \in \R,\; \bm{x}_i \in \mathcal{X}\right\}
\end{equation}
muni du produit scalaire~:
\begin{equation}
  \left\langle \sum_{i} \alpha_i\, \kappa(\bm{x}_i, \cdot),\;
  \sum_{j} \beta_j\, \kappa(\bm{x}_j', \cdot)
  \right\rangle_{\mathcal{H}_\kappa}
  = \sum_{i,j} \alpha_i \beta_j\, \kappa(\bm{x}_i, \bm{x}_j')
\end{equation}
\end{proposition}

% ============================================================
\section{Th\'eor\`eme du repr\'esentant}
\label{sec:representer}
\index{th\'eor\`eme du repr\'esentant}

\begin{theoreme}[Th\'eor\`eme du repr\'esentant --- Kimeldorf \& Wahba, 1971]
\label{thm:representer}
Soit $\Omega : \R_+ \to \R$ une fonction croissante et
$L : (\mathcal{X} \times \R \times \R)^n \to \R \cup \{+\infty\}$ une
fonction de perte quelconque. Le probl\`eme d'optimisation~:
\begin{equation}
  \min_{f \in \mathcal{H}_\kappa} L\!\left((\bm{x}_1, y_1, f(\bm{x}_1)),
  \ldots, (\bm{x}_n, y_n, f(\bm{x}_n))\right)
  + \Omega(\norm{f}^2_{\mathcal{H}_\kappa})
  \label{eq:representer-problem}
\end{equation}
admet un minimiseur de la forme~:
\begin{equation}
  f^*(\bm{x}) = \sum_{i=1}^n \alpha_i\, \kappa(\bm{x}_i, \bm{x})
  \label{eq:representer-solution}
\end{equation}
pour certains coefficients $\alpha_1, \ldots, \alpha_n \in \R$.
\end{theoreme}

\begin{proof}[Esquisse de preuve]
D\'ecomposons $f \in \mathcal{H}_\kappa$ en une partie dans le
sous-espace engendr\'e par $\{\kappa(\bm{x}_i, \cdot)\}_{i=1}^n$
et son compl\'ement orthogonal~:
\begin{equation}
  f = \underbrace{\sum_{i=1}^n \alpha_i\, \kappa(\bm{x}_i, \cdot)}_{f_{\parallel}}
  + \underbrace{f_\perp}_{f_\perp \perp \text{span}\{\kappa(\bm{x}_i,\cdot)\}}
\end{equation}
Par la propri\'et\'e de reproduction,
$f(\bm{x}_j) = f_{\parallel}(\bm{x}_j) + \inner{f_\perp}{\kappa(\bm{x}_j, \cdot)} = f_{\parallel}(\bm{x}_j)$,
donc $f_\perp$ n'affecte pas le terme de perte $L$. Par ailleurs~:
\begin{equation}
  \norm{f}^2_{\mathcal{H}_\kappa} = \norm{f_{\parallel}}^2 +
  \norm{f_\perp}^2 \geq \norm{f_{\parallel}}^2
\end{equation}
Puisque $\Omega$ est croissante, ajouter $f_\perp \neq 0$ augmente la
p\'enalit\'e sans diminuer la perte. Donc le minimiseur v\'erifie
$f_\perp = 0$.
\end{proof}

\begin{intuition}[title=\textbf{Port\'ee du th\'eor\`eme du repr\'esentant}]
Ce th\'eor\`eme est fondamental car il transforme un probl\`eme
d'optimisation en \textbf{dimension infinie} (minimiser sur
$\mathcal{H}_\kappa$) en un probl\`eme en \textbf{dimension finie}
(trouver $\bm{\alpha} \in \R^n$). C'est ce qui rend les m\'ethodes
\`a noyaux calculatoirement r\'ealisables.
\end{intuition}

% ============================================================
\section{Connexions avec les processus gaussiens}
\label{sec:gp-connection}
\index{processus gaussiens}

\begin{definition}[Processus gaussien]
Un \textbf{processus gaussien} (GP) est une collection de variables
al\'eatoires dont toute sous-collection finie suit une distribution
gaussienne multivari\'ee. Un GP est enti\`erement d\'efini par sa
fonction moyenne $\mu(\bm{x}) = \E[f(\bm{x})]$ et sa fonction de
covariance $\kappa(\bm{x}, \bm{x}') = \Cov[f(\bm{x}), f(\bm{x}')]$.
\end{definition}

\begin{proposition}[Lien GP / Kernel Ridge Regression]
\label{prop:gp-krr}
La moyenne a posteriori du GP avec prior $f \sim \mathcal{GP}(0, \kappa)$
et vraisemblance gaussienne $y_i | f \sim \mathcal{N}(f(\bm{x}_i),
\sigma^2)$ co\"incide avec la solution de la r\'egression Ridge \`a
noyau~:
\begin{equation}
  \E[f(\bm{x}) | \bm{y}] = \bm{k}(\bm{x})^\top
  (\bm{K} + \sigma^2 \bm{I})^{-1} \bm{y}
  \label{eq:gp-posterior-mean}
\end{equation}
o\`u $k(\bm{x})_i = \kappa(\bm{x}_i, \bm{x})$ et $\lambda = \sigma^2 / n$
dans~\eqref{eq:krr-alpha}.
\end{proposition}

\begin{remarque}
Le GP fournit en plus une \textbf{incertitude pr\'edictive}~:
\begin{equation}
  \Var[f(\bm{x}) | \bm{y}] = \kappa(\bm{x}, \bm{x}) -
  \bm{k}(\bm{x})^\top (\bm{K} + \sigma^2 \bm{I})^{-1} \bm{k}(\bm{x})
\end{equation}
C'est un avantage majeur des GP par rapport \`a la Kernel Ridge
Regression~: on obtient non seulement une pr\'ediction ponctuelle,
mais aussi une barre d'erreur.
\end{remarque}

% ============================================================
\section{Impl\'ementation~: m\'ethodes \`a noyaux avec scikit-learn}
\label{sec:implementation-kernels}

\begin{codeblock}[title=\textbf{SVM \`a noyau, Kernel Ridge et Kernel PCA}]
\begin{minted}{python}
import numpy as np
import matplotlib.pyplot as plt
from sklearn.datasets import make_moons, make_circles
from sklearn.svm import SVC
from sklearn.kernel_ridge import KernelRidge
from sklearn.decomposition import KernelPCA
from sklearn.preprocessing import StandardScaler
from sklearn.model_selection import GridSearchCV
from sklearn.pipeline import Pipeline

# --- Donnees non lineairement separables ---
X, y = make_moons(n_samples=300, noise=0.2, random_state=42)
scaler = StandardScaler()
X_scaled = scaler.fit_transform(X)

# --- SVM avec differents noyaux ---
fig, axes = plt.subplots(1, 3, figsize=(15, 4))
kernels = ['linear', 'poly', 'rbf']
for ax, kernel in zip(axes, kernels):
    svm = SVC(kernel=kernel, degree=3, gamma='scale', C=1.0)
    svm.fit(X_scaled, y)
    # Grille de decision
    xx, yy = np.meshgrid(
        np.linspace(X_scaled[:,0].min()-0.5,
                    X_scaled[:,0].max()+0.5, 200),
        np.linspace(X_scaled[:,1].min()-0.5,
                    X_scaled[:,1].max()+0.5, 200)
    )
    Z = svm.predict(np.c_[xx.ravel(), yy.ravel()])
    Z = Z.reshape(xx.shape)
    ax.contourf(xx, yy, Z, alpha=0.3, cmap='coolwarm')
    ax.scatter(X_scaled[:,0], X_scaled[:,1], c=y,
               cmap='coolwarm', edgecolors='k', s=20)
    acc = svm.score(X_scaled, y)
    ax.set_title(f"{kernel} (acc={acc:.2f})")
plt.tight_layout()
plt.savefig("svm_kernels.pdf")
plt.show()
\end{minted}
\end{codeblock}

\begin{codeblock}[title=\textbf{Kernel PCA~: r\'eduction non lin\'eaire}]
\begin{minted}{python}
# --- Kernel PCA sur donnees circulaires ---
X_circ, y_circ = make_circles(n_samples=400, factor=0.3,
                               noise=0.1, random_state=42)

fig, axes = plt.subplots(1, 3, figsize=(15, 4))

# PCA lineaire
from sklearn.decomposition import PCA
pca = PCA(n_components=2)
X_pca = pca.fit_transform(X_circ)
axes[0].scatter(X_pca[:,0], X_pca[:,1], c=y_circ,
                cmap='coolwarm', s=15)
axes[0].set_title("PCA lineaire")

# Kernel PCA avec noyau RBF
kpca_rbf = KernelPCA(n_components=2, kernel='rbf', gamma=5)
X_kpca_rbf = kpca_rbf.fit_transform(X_circ)
axes[1].scatter(X_kpca_rbf[:,0], X_kpca_rbf[:,1],
                c=y_circ, cmap='coolwarm', s=15)
axes[1].set_title("Kernel PCA (RBF, $\\gamma=5$)")

# Kernel PCA avec noyau polynomial
kpca_poly = KernelPCA(n_components=2, kernel='poly',
                       degree=3, gamma=1)
X_kpca_poly = kpca_poly.fit_transform(X_circ)
axes[2].scatter(X_kpca_poly[:,0], X_kpca_poly[:,1],
                c=y_circ, cmap='coolwarm', s=15)
axes[2].set_title("Kernel PCA (poly, $p=3$)")

plt.tight_layout()
plt.savefig("kernel_pca.pdf")
plt.show()
\end{minted}
\end{codeblock}

\begin{codeblock}[title=\textbf{Kernel Ridge Regression}]
\begin{minted}{python}
# --- Kernel Ridge Regression sur donnees 1D ---
np.random.seed(42)
X_train = np.sort(np.random.uniform(-3, 3, 50))[:, None]
y_train = np.sin(X_train).ravel() + 0.3 * np.random.randn(50)
X_plot = np.linspace(-3.5, 3.5, 300)[:, None]

fig, axes = plt.subplots(1, 3, figsize=(15, 4))
gammas = [0.5, 2.0, 10.0]

for ax, gamma in zip(axes, gammas):
    krr = KernelRidge(kernel='rbf', alpha=0.1, gamma=gamma)
    krr.fit(X_train, y_train)
    y_pred = krr.predict(X_plot)

    ax.scatter(X_train, y_train, c='steelblue', s=20,
               label="Donnees")
    ax.plot(X_plot, np.sin(X_plot), 'g--', label="$\\sin(x)$")
    ax.plot(X_plot, y_pred, 'r-', lw=2,
            label=f"KRR ($\\gamma={gamma}$)")
    ax.set_title(f"$\\gamma = {gamma}$")
    ax.legend(fontsize=8)
    ax.set_ylim(-2, 2)

plt.tight_layout()
plt.savefig("kernel_ridge.pdf")
plt.show()
\end{minted}
\end{codeblock}

\begin{codeoutput}
\begin{verbatim}
[Figures sauvegardees: svm_kernels.pdf, kernel_pca.pdf,
 kernel_ridge.pdf]
\end{verbatim}
\end{codeoutput}

% ============================================================
\section{Exercices}
\label{sec:exercises-ch12}

\begin{exercice}[Feature maps explicites --- niveau 1]
\begin{enumerate}
  \item Pour le noyau polynomial $\kappa(\bm{x}, \bm{x}') =
    (1 + \bm{x}^\top \bm{x}')^2$ avec $\bm{x} \in \R^2$, d\'eterminer
    explicitement la feature map $\phi : \R^2 \to \R^p$.
    Quelle est la valeur de $p$~?
  \item V\'erifier que $\kappa(\bm{x}, \bm{x}') =
    \inner{\phi(\bm{x})}{\phi(\bm{x}')}$ sur l'exemple num\'erique
    $\bm{x} = (1, 2)^\top$, $\bm{x}' = (3, -1)^\top$.
  \item La fonction $\kappa(\bm{x}, \bm{x}') = (\bm{x}^\top \bm{x}')^2
    - 1$ est-elle un noyau valide~? Justifier.
\end{enumerate}
\end{exercice}

\begin{exercice}[Th\'eor\`eme du repr\'esentant et Kernel Ridge --- niveau 2]
\begin{enumerate}
  \item Montrer que le probl\`eme de r\'egression Ridge \`a noyau
    \eqref{eq:krr-functional}, par application du th\'eor\`eme du
    repr\'esentant, se r\'eduit \`a~:
    \begin{equation*}
      \min_{\bm{\alpha} \in \R^n} \frac{1}{n}
      \norm{\bm{y} - \bm{K}\bm{\alpha}}^2 +
      \lambda\, \bm{\alpha}^\top \bm{K} \bm{\alpha}
    \end{equation*}
  \item En d\'erivant par rapport \`a $\bm{\alpha}$ et en annulant le
    gradient, retrouver la solution~\eqref{eq:krr-alpha}.
  \item Impl\'ementer la Kernel Ridge Regression \`a partir de z\'ero
    (sans \texttt{KernelRidge}) avec un noyau RBF. Comparer vos
    pr\'edictions avec celles de scikit-learn.
\end{enumerate}
\end{exercice}

\begin{exercice}[Kernel PCA et processus gaussiens --- niveau 3]
\begin{enumerate}
  \item G\'en\'erer un jeu de donn\'ees en forme de spirale
    ($n = 500$ points, 2 classes). Appliquer la Kernel PCA avec
    diff\'erents noyaux et valeurs de $\gamma$.
    Visualiser les r\'esultats et identifier la meilleure combinaison.
  \item Impl\'ementer le $k$-means \`a noyau \`a partir de z\'ero en
    utilisant uniquement la matrice de Gram. Comparer avec
    \texttt{SpectralClustering} de scikit-learn.
  \item Utiliser \texttt{GaussianProcessRegressor} de scikit-learn pour
    r\'egresser $y = \sin(x) + \varepsilon$ avec un noyau RBF.
    Tracer la moyenne a posteriori et l'intervalle de confiance \`a
    $\pm 2\sigma$. V\'erifier num\'eriquement que la moyenne a
    posteriori co\"incide avec la pr\'ediction de
    \texttt{KernelRidge} (avec $\alpha = \sigma^2_{\text{bruit}}$).
  \item (\textbf{Th\'eorique}) Montrer que la somme de deux noyaux
    d\'efinis positifs est un noyau d\'efini positif. Le produit de
    deux noyaux d\'efinis positifs est-il aussi un noyau d\'efini
    positif~? Prouver ou donner un contre-exemple.
\end{enumerate}
\end{exercice}
